\section{Conclusions and future work}
\label{sec:conclusion}

\subsection{\PurePy-compatible standard libraries}

Scientific computing in Python rests on libraries such as NumPy and Pandas, whose interfaces lean on
mutation and indexing. \PurePy therefore needs compatible counterparts: libraries, or subsets of existing
ones, whose operations are pure. Polars,\footnote{\url{https://docs.pola.rs/}} whose operations return new frames rather than
updating in place, suggests the shape such interfaces can take.

\subsection{Typed effects}

A future version of \PurePy may add a typed effect discipline. In a mypy-compatible type system the
question is notational, because effect information must fit within annotations that Python and mypy
already accept. A return type such as \pyinline{IO[int]} tracks effects as an ordinary generic type, in
the Haskell style, although chaining is awkward in a language of statements without dedicated syntax.
\pyinline{Annotated[int, Console]} (PEP 593) attaches effect metadata that mypy ignores and the \PurePy
checker can read. Lightest is a decorator: with effects the default, \pyinline{@pure} marks functions in
which printing and other effects are rejected, and unannotated code keeps its current meaning. Under any
of these, existing support for printing is not an obstacle, since printing can be reinterpreted as a typed
output effect or kept native with essentially equivalent semantics, printed output never influencing a
result (\secref{purity}). Typed disciplines coexisting with native effects have precedent, from Haskell's
untracked \pyinline{trace} in pure code\footnote{\url{https://hackage.haskell.org/package/base/docs/Debug-Trace.html}}
to Scala's monadic effect libraries alongside the language's native effects.

\subsection{Proof of agreement with Python}
\label{sec:agreement}

The guarantee of \secref{python-compatibility} is informal. To prove it we would need a reference semantics
of Python to state it against; one candidate is the core calculus of Politz et al.~\cite{politz2013}, which
covers a subset of Python and predates the current language.
